% Part: second-order-logic
% Chapter: syntax-and-semantics
% Section: terms-formulas

\documentclass[../../../include/open-logic-section]{subfiles}

\begin{document}

\olfileid{sol}{syn}{frm}

\olsection{পদ ও \printtoken{P}{formula}}

প্রথম-ক্রমের যুক্তিবিদ্যার মতোই দ্বিতীয়-ক্রমের যুক্তিবিদ্যার রাশিগুলি
\emph{!!{variable}s}, \emph{!!{constant}s}, \emph{!!{predicate}s} এবং
কখনও \emph{!!{function}s}-সম্বলিত একটি মৌলিক শব্দভাণ্ডার থেকে গঠিত হয়।
এগুলির সঙ্গে যৌক্তিক সংযোজক, পরিমাণসূচক এবং বন্ধনী ও কমার মতো যতিচিহ্ন
মিলিয়ে \emph{পদ} ও \emph{!!{formula}s} তৈরি হয়। পার্থক্য হলো,
দ্বিতীয়-ক্রমের যুক্তিবিদ্যায় বস্তু-চলরাশির পাশাপাশি সম্পর্ক ও অপেক্ষকের
জন্যও চলরাশি থাকে এবং তাদের উপর পরিমাণায়ন করা যায়। তাই প্রথম-ক্রমের
যুক্তিবিদ্যার যৌক্তিক সংকেতগুলির সঙ্গে দ্বিতীয়-ক্রমের যুক্তিবিদ্যায় আরও
থাকে:

\begin{enumerate}
\item প্রত্যেক স্থানসংখ্যা~$n$-এর দ্বিতীয়-ক্রমীয় সম্পর্ক-!!{variable}s-এর
  একটি !!{denumerable}s সেট: $\Obj V_0^n$, $\Obj V_1^n$, $\Obj V_2^n$, \dots
\item দ্বিতীয়-ক্রমীয় অপেক্ষক-!!{variable}s-এর একটি !!{denumerable}s সেট:
  $\Obj u_0^n$, $\Obj u_1^n$, $\Obj u_2^n$, \dots
\end{enumerate}

প্রথম-ক্রমীয় চলরাশি $\Obj v_i$-এর অধিচলরাশি হিসেবে আমরা যেমন $x$,
$y$, $z$ ব্যবহার করি, তেমনি $\Obj V_i^n$-এর অধিচলরাশি হিসেবে $X$,
$Y$, $Z$ ইত্যাদি এবং~$\Obj u_i^n$-এর অধিচলরাশি হিসেবে $u$, $v$
ইত্যাদি ব্যবহার করব।

\begin{explain}
দ্বিতীয়-ক্রমীয় কোনো ভাষার অ-যৌক্তিক সংকেতগুলি প্রথম-ক্রমীয় ভাষার মতোই
নির্দিষ্ট করা হয়: তার !!{constant}s, !!{function}s ও !!{predicate}s-এর
তালিকা দিয়ে।

প্রথম-ক্রমের যুক্তিবিদ্যায় সাধারণত !!{identity}~$\eq$ অন্তর্ভুক্ত থাকে।
প্রথম-ক্রমের যুক্তিবিদ্যায় কোনো ভাষা~$\Lang{L}$-এর অ-যৌক্তিক সংকেতগুলি
আকর্ষণীয় কিছু প্রকাশের জন্য অত্যন্ত জরুরি। অ-যৌক্তিক সংকেতবিহীন
!!{sentence}s অবশ্যই আছে, কিন্তু শুধু~$\eq$ দিয়ে আকর্ষণীয় কিছু বলা
কঠিন। দ্বিতীয়-ক্রমের যুক্তিবিদ্যায় সম্পর্ক ও অপেক্ষক-চলরাশির অফুরন্ত
ভাণ্ডার আছে বলে অ-যৌক্তিক সংকেতের বিশেষ কোনো ভাণ্ডার ছাড়াই আমরা
প্রথম-ক্রমীয় ভাষায় বলা যায় এমন সব কথা বলতে পারি।
\end{explain}

\begin{defn}[দ্বিতীয়-ক্রমীয় পদ]
~$\Lang L$-এর \emph{দ্বিতীয়-ক্রমীয় পদ}-এর সেট $\TrmSOL[L]$,
\olref[fol][syn][frm]{defn:terms}-এর সঙ্গে নিচের দফাটি যোগ করে সংজ্ঞায়িত
হয়:
\begin{enumerate}
\item $u$ যদি একটি $n$-স্থানীয় অপেক্ষক-চলরাশি হয় এবং $t_1$, \dots,
  $t_n$ যদি পদ হয়, তবে $\Atom{u}{t_1, \ldots, t_n}$ একটি পদ।
\end{enumerate}
\end{defn}

\begin{explain}
সুতরাং, দ্বিতীয়-ক্রমীয় পদ দেখতে প্রথম-ক্রমীয় পদের মতোই; শুধু
প্রথম-ক্রমীয় পদে যেখানে !!a{function}~$\Obj{f^n_i}$ থাকে,
দ্বিতীয়-ক্রমীয় পদে তার জায়গায় অপেক্ষক-চলরাশি~$\Obj{u^n_i}$ থাকতে
পারে।
\end{explain}

\begin{defn}[দ্বিতীয়-ক্রমীয় \usetoken{s}{formula}]
ভাষা~$\Lang L$-এর \emph{দ্বিতীয়-ক্রমীয় !!{formula}s}-এর সেট
~$\FrmSOL[L]$, \olref[fol][syn][frm]{defn:formulas}-এর সঙ্গে নিচের
দফাগুলি যোগ করে সংজ্ঞায়িত হয়:
\begin{enumerate}
\item $X$ যদি একটি $n$-স্থানীয় বিধেয়-চলরাশি হয় এবং $t_1$, \dots,
  $t_n$ যদি~$\Lang L$-এর দ্বিতীয়-ক্রমীয় পদ হয়, তবে
  $\Atom{X}{t_1,\ldots, t_n}$ একটি পরমাণু !!{formula}।

\tagitem{prvAll}{$!A$ যদি !!a{formula} এবং $u$ যদি একটি অপেক্ষক-চলরাশি
  হয়, তবে $\lforall[u][!A]$ একটি !!a{formula}।}{}

\tagitem{prvAll}{$!A$ যদি !!a{formula} এবং $X$ যদি একটি বিধেয়-চলরাশি
  হয়, তবে $\lforall[X][!A]$ একটি !!a{formula}।}{}

\tagitem{prvEx}{$!A$ যদি !!a{formula} এবং $u$ যদি একটি অপেক্ষক-চলরাশি
  হয়, তবে $\lexists[u][!A]$ একটি !!a{formula}।}{}

\tagitem{prvEx}{$!A$ যদি !!a{formula} এবং $X$ যদি একটি বিধেয়-চলরাশি
  হয়, তবে $\lexists[X][!A]$ একটি !!a{formula}।}{}
\end{enumerate}
\end{defn}

\end{document}
